\section{Robustness and Institutional Boundary Calculations}
\label{app:robustness}
This appendix supplies the proof bridges for the scope statements in Section 5
of the short article. Each alternative is fixed before the comparison and
re-solved at its own primitives. These are specification boundaries, not
theorems about all network or differentiated-demand models.

\subsection{Canonical network strength}
Use the canonical blocks in Section 2 and define
\[
\Phi(c,v)=K_M+2A+C-K_I-3P,
\qquad N(c,v)=32(1-v)^2(2-3v)^2\Phi(c,v).
\]
Thus $N$ is a polynomial of degree four in $v$, obtained from the displayed
Cournot solutions, and its sign is the government's SU-minus-IS ranking:
the cleared denominator is strictly positive. At $v=0$,
\[
\Phi(c,0)=\frac{c(13c-6)}{32}<0\quad(0<c<1/3).
\]
The adoption thresholds nevertheless obey
\begin{align*}
T_A(0)&=\frac{3c(2-3c)}{16},&
T_U(0)=T_W(0)&=\frac{3c(2-c)}{16},\\
2T_A(0)-T_W(0)&=\frac{3c(2-5c)}{16}>0.
\end{align*}
Consequently one-way adoption can occur without a favorable initial SU ranking.

For completeness, the fixed-$c$ network comparison has a precise root target.
Feasibility is $0<v<U(c)$, where
$U(c)=\min\{1/4,(1-3c)/[3(1-c)]\}$. The change
$v=U(c)t/(1+t)$ maps this open interval to $t>0$.
For $0<c\le1/9$, the polynomial
\[
L_1(t;c)=(1+t)^4N\left(c,\frac{t}{4(1+t)}\right)
\]
has leading coefficient $9(145c^2-18c+3)/64>0$ and four remaining
coefficients strictly negative; its constant is $4c(13c-6)$.
For $1/9<c<1/3$, use
\[
L_2(t;c)=9(1-c)^3(1+t)^4
N\left(c,\frac{(1-3c)t}{3(1-c)(1+t)}\right).
\]
Its leading coefficient is $Q(c)=17c^3+109c^2-89c+11$ and all
remaining coefficients are strictly negative. These coefficient claims are
exact univariate sign checks: the coefficients of $t^j$, $j=1,2,3$, have no
root in their stated $c$ interval and are negative at $c=1/18$ for $L_1$
and $c=1/8$ for $L_2$. The endpoint $c=1/9$ uses $L_1$.
Root counts use Sturm's sequence $s_0=a$, $s_1=a'$,
$s_{k+1}=-\operatorname{rem}(s_{k-1},s_k)$; the difference in sign
variations at rational endpoints, after omitting zeros, counts roots.
The constants are $4c(13c-6)$ for $L_1$ and
$36c(1-c)^3(13c-6)$ for $L_2$, strictly negative for $0<c<1/3$;
the excluded endpoint $c=0$ is not used as a negative test value.

There is exactly one root $c^*$ of $Q$ in $(0,1/3)$, with opposite signs
at the rational endpoints of its isolating interval
\[
c^*\in(1529167/10^7,1529168/10^7).
\]
Descartes' rule therefore gives exactly one feasible network threshold when
$0<c<c^*$: $\Phi<0$ below it and $\Phi>0$ above it. When
$c\ge c^*$, all nonzero coefficients of $L_2$ are negative and no
initial SU advantage exists. The excluded upper endpoint $v=U(c)$ is not an
interior equilibrium claim. The mechanical coefficient/root-count checks are
in \path{code/revision/check_r4_network_role.py}; the mathematical target is
$N$, not a numerical grid.

\subsection{Exact polynomial sign certificates}
The certificates below use an additional exact route, independent of the
canonical root calculation. For a polynomial
$f(x,y)=\sum_{k=0}^{n}\sum_{\ell=0}^{m}a_{k\ell}x^ky^\ell$, write
\begin{align*}
f(x,y)&=\sum_{i=0}^{n}\sum_{j=0}^{m}
 b_{ij}\binom ni x^i(1-x)^{n-i}\binom mj y^j(1-y)^{m-j},\\
b_{ij}&=\sum_{k\le i,\,\ell\le j}a_{k\ell}
 \frac{\binom ik}{\binom nk}\frac{\binom j\ell}{\binom m\ell}.
\end{align*}
All operations are rational. On $[0,1]^2$ the basis weights are nonnegative
and sum to one. Nonnegative coefficients with at least one positive active
weight prove positivity. The reported minimum positive coefficient and the
complete zero-index set specify the certificate; the conversion formula and
the polynomial definitions specify its input. Long coefficient lists may be
regenerated, but neither the object nor the inference is delegated to code.

\subsection{Network benefits for singleton groups}
The alternative rule is $g_i=v\sum_{j\in G_i}q_j$ for every nonempty
group, including singletons. Inverse demand remains $p_i=1-\sum_jq_j+g_i$.
The own-output derivative gives operating profit $(1-v)q_i^2$ at an
interior Cournot equilibrium. In a member market the two equal member outputs
$x$ and outsider output $y$ solve
\[
1-3(1-v)x-y=0,\qquad 1-c-2x-2(1-v)y=0.
\]
Put $K=3v^2-6v+2$. The member and outsider markets yield, respectively,
\begin{align*}
q_M^{S1}&=\frac{1+c-2v}{2K},&
q_B^{S1}&=\frac{1-3v-3c(1-v)}{2K},\\
q_C^{S1}&=\frac{1-2v-2c(1-v)}{2K},&
q_D^{S1}&=\frac{1+2c-3v}{2K}.
\end{align*}
The complete-compatibility IS blocks remain $P=1/[16(1-v)]$ and
$K_I=9/[32(1-v)^2]$. Define explicitly
\begin{align*}
A_1&=(1-v)(q_M^{S1})^2,& C_1&=(1-v)(q_C^{S1})^2,\\
K_{M1}&=(2q_M^{S1}+q_B^{S1})^2/2,&
\Phi_{S1}&=K_{M1}+2A_1+C_1-K_I-3P.
\end{align*}
On the canonical domain $0<c<c_{\max}(v)=(1-3v)/[3(1-v)]$,
$0\le v<1/4$, $K\ge K(1/4)=11/16>0$. The outsider numerator is
the feasibility slack; the $q_C^{S1}$ numerator exceeds its value $1/3$
at $c=c_{\max}(v)$, and the other two are positive directly.
Thus these are interior, uniquely solved Cournot configurations.

Let $M_1=32(1-v)^2K^2\Phi_{S1}$ and
\[
Q_1(x,y)=-9M_1\left(xc_{\max}(y/4),y/4\right).
\]
Cancellation makes $Q_1$ a polynomial of degree $(2,5)$. Its exact Bernstein
coefficient array (rows $i=0,1,2$; columns $j=0,\ldots,5$) is
\[
\begin{pmatrix}
0&54/5&63/5&1611/160&8019/1280&1539/512\\
36&138/5&387/20&1947/160&8439/1280&1467/512\\
20&74/5&201/20&979/160&4023/1280&571/512
\end{pmatrix}.
\]
Here $0<x<1$ and $0\le y<1$. Positive weights occur even at $y=0$
because $x>0$. Hence $Q_1>0$, $M_1<0$, and $\Phi_{S1}<0$ throughout
the claimed domain. The equality at $(c,v)=(0,0)$ is excluded.
This proves a specification failure of the initial regional ranking.

Private adoption need not fail. Under $SW$ the foreign singleton quantity is
\[
q_S^{S1}=\frac{1-2v-2c(1-v)}{2(2-v)(1-2v)}.
\]
Set $B_1=(1-v)(q_B^{S1})^2$, $S_1=(1-v)(q_S^{S1})^2$ and
$(T_A^{S1},T_U^{S1},T_W^{S1})=(P-B_1,A_1-C_1,A_1-S_1)$.
At $(c,v)=(1/10,6/25)$ their exact values are
\[
\left(\frac{5183091}{63768256},
\frac{295659}{3356224},
\frac{7226119311}{68631424576}\right),
\]
and their maximum is less than $2T_A^{S1}$. Against the two possible
rival actions the reverse gains are $T_U^{S1}$ and $T_A^{S1}$; the $SW$
gains are $T_W^{S1}$ and $T_A^{S1}$. Thus choosing $F$ strictly between
that maximum and $2T_A^{S1}$ selects outsider-only adoption.

\subsection{Differentiated Cournot and Bertrand demand}
Fix $\gamma=1/2$. Let $C_G$ have entry one for two firms in the same
compatibility group of size at least two, including that group's diagonal,
and zero otherwise. In a market with cost vector $r$, set
\[
M_G=\tfrac12 I+\tfrac12\mathbf1\mathbf1^\top-vC_G,
\qquad p=\mathbf1-M_Gq,
\qquad CS=\tfrac12q^\top M_Gq.
\]
The four configurations needed for all blocks are
\begin{center}\small
\begin{tabular}{@{}lll@{}}\toprule
Configuration & Compatibility groups & Cost vector $r$\\\midrule
$I$ (IS) & $\{1,2,3\}$ & $(0,0,0)$\\
$M$ (SU member market) & $\{1,2\},\{3\}$ & $(0,0,c)$\\
$O$ (SU outsider market) & $\{1,2\},\{3\}$ & $(c,c,0)$\\
$W$ (SW market) & $\{1\},\{2\},\{3\}$ & $(0,c,c)$\\\bottomrule
\end{tabular}\end{center}
The audited domain is the full box $0<c<1/3$, $0\le v<1/4$.
For $I,M,O$, the first two leading principal minors are $1-v$ and
$3/4-v$; the determinants are $(2-3v)/4$ for $I$ and
$(1-2v)/2$ for $M,O$. For $W$, the matrix has eigenvalues
$2,1/2,1/2$. All are positive, so demand is positive definite.

\paragraph{Quantity competition.}
Firm $i$ maximizes $(p_i-r_i)q_i$. Its first-order condition and solution are
\[
1-r_i-2M_{ii}q_i-\sum_{j\ne i}M_{ij}q_j=0,
\qquad q^{C}=(M+\operatorname{diag}(M_{ii}))^{-1}(\mathbf1-r).
\]
Own profit curvature is $-2M_{ii}<0$, and the coefficient matrix is positive
definite. The outputs in configurations $I,M,O,W$ are
\begin{align*}
q^C_I&=\frac{\mathbf1}{3-4v},\\
q^C_M&=\frac{(c+3,c+3,3-6v-c(5-6v))^\top}{3(3-4v)},\\
q^C_O&=\frac{(3-4c,3-4c,3+2c-6v)^\top}{3(3-4v)},\\
q^C_W&=\frac{(3+2c,3-4c,3-4c)^\top}{9}.
\end{align*}
All outputs are positive: the potentially binding $q^C_M$ numerator exceeds
$4/3-4v>0$, its value at $c=1/3$; the other bounds are immediate.
Equilibrium operating profit is $M_{ii}(q_i^C)^2$.

\paragraph{Price competition.}
Direct demand is $q=D_G(\mathbf1-p)$ with $D_G=M_G^{-1}$, not the
canonical unit-demand formula. Maximizing $(p_i-r_i)q_i$ yields
\[
q_i-D_{ii}(p_i-r_i)=0,\qquad
p^B=(D+\operatorname{diag}(D_{ii}))^{-1}
 [D\mathbf1+\operatorname{diag}(D_{ii})r].
\]
Own profit curvature is $-2D_{ii}<0$. Positive definiteness gives a unique
solution of these interior price conditions. Write $h=8v^2-16v+7$,
$\ell=(1-2v)h$, $a=3-4v$, and
\begin{align*}
b_M&=7-12v+c(3-4v),& b_C&=7-12v+c(16v-10),\\
b_B&=16v^2-24v+7-c(16v^2-32v+13),\\
b_D&=16v^2-24v+7+c(6-8v).
\end{align*}
The resulting quantities are
\begin{align*}
q_I^B&=\frac{a\mathbf1}{4(1-v)(2-3v)},&
q_W^B&=\frac{3(7+6c,7-10c,7-10c)^\top}{56},\\
q_M^B&=\frac{a(b_M,b_M,b_B)^\top}{8\ell},&
q_O^B&=\frac{a(b_C,b_C,b_D)^\top}{8\ell}.
\end{align*}
Here $h\ge h(1/4)=7/2>0$, so every displayed denominator is positive.
The lower bound for $b_B$ at $c=1/3$ is
$\tfrac83(4v-1)(v-1)>0$; $b_C>2$, $b_D>2$, and $b_M>4$.
All quantities are therefore interior on the stated box. Operating profit is
$(q_i^B)^2/D_{ii}$, obtained from the price first-order condition.

\paragraph{Welfare targets and strict signs.}
For each mode $X\in\{C,B\}$, obtain $P^X$ from firm 1 in $I$,
$A^X,B^X$ from firms 1 and 3 in $M$, $C^X,D^X$ from firms 1 and 3
in $O$, and $S^X$ from firm 2 in $W$. Let $K_I^X,K_M^X$ be the
corresponding quadratic consumer surpluses. National welfare then gives
\begin{align*}
\Phi^X&=K_M^X+2A^X+C^X-K_I^X-3P^X,\\
\mathscr E^X&=K_M^X-K_I^X+2(A^X-P^X),&
\mathscr D^X&=P^X-C^X.
\end{align*}
In quantity competition the two sign numerators are explicitly
\begin{align*}
N_C&=2(18v^2-44v+27)c^2-12(6v^2-10v+5)c+9v(4v-3),\\
E_C&=2(18v^2-28v+11)c^2-12(6v^2-6v+1)c+9v(4v-3),
\end{align*}
with $(\Phi^C,\mathscr E^C)=(N_C,E_C)/[18(3-4v)^2]$ and
$\mathscr D^C=8c(3-2c)(1-v)/[9(3-4v)^2]>0$.

For price competition define the strictly positive denominator
$D_+=(1-v)^2(1-2v)(2-3v)h^2$ and the polynomial targets
\[
Q_\Phi=\frac{64D_+\Phi^B}{4v-3},\qquad
Q_E=\frac{64D_+\mathscr E^B}{4v-3},\qquad
Q_D=\frac{32D_+\mathscr D^B}{3-4v}.
\]
These definitions, together with the explicit equilibrium and profit formulas,
identify the polynomials without printing their long expansions. Apply the
Bernstein conversion above with $(c,v)=(x/3,y/4)$. Exact certificates are
\begin{center}\small
\begin{tabular}{@{}llll@{}}\toprule
Polynomial & Degree $(n,m)$ & Zero indices & Minimum positive $b_{ij}$\\\midrule
$-N_C$ & $(2,2)$ & $(0,0)$ & $27/8$\\
$-E_C$ & $(2,2)$ & $(0,0)$ & $14/9$\\
$Q_\Phi$ & $(2,7)$ & $(0,0)$ & $51/8$\\
$Q_E$ & $(2,7)$ & $(0,0),(2,7)$ & $11/6$\\
$Q_D$ & $(2,5)$ & $(0,0)$ & $1$\\\bottomrule
\end{tabular}\end{center}
Every unlisted coefficient is positive. For $0<x<1$, $0<y<1$ the
positive weights prove strict positivity; at $y=0$ the positive coefficients
with $i>0,j=0$ do so. The corner $(x,y)=(1,1)$ is outside the audited
domain, and equality at $(0,0)$ does not weaken any claim with $c>0$.
Since $4v-3<0$, the signs are
$\Phi^C,\Phi^B,\mathscr E^C,\mathscr E^B<0$ and
$\mathscr D^C,\mathscr D^B>0$ throughout the claimed box.

For full-cost-relief outsider adoption the member market becomes $I$ but the
outsider market remains $O$, giving the bridge
\[
W_M^{SU,O,X}-W^{IS,X}=-\mathscr D^X<0,\qquad
\Delta(W^{IS,X}-W_M^{SU,X})=\mathscr E^X<0.
\]
The member prefers IS both before and after adoption; adoption weakens its
relative IS incentive. With $T_A^X=P^X-B^X$, $T_U^X=A^X-C^X$ and
$T_W^X=A^X-S^X$, exact substitution at $(c,v)=(1/10,6/25)$ gives
\[
\max\{T_A^X,T_U^X,T_W^X\}<1/5<2T_A^X
\quad(X=C,B).
\]
The same two reverse/$SW$ deviation cases as above show that $F=1/5$
selects outsider-only adoption. Thus the political failure is compatible
with privately optimal adaptation. It occurs under both competition modes
at the same demand primitives, so it cannot be attributed to Bertrand alone.
No claim is made for other $\gamma$, network rules, or residual-cost
coalition games.

\paragraph{Verification artifacts.}
The original R4/R5 checks are retained in the replication package.
\path{code/revision/check_jict_short_independent.py} separately reconstructs
the matrix equilibria, the S1 triangular-domain certificate, the displayed
Bernstein certificates, and every firm action at the exact witnesses. Its
optional JSON output includes all rational coefficients and a reconstruction
of each polynomial from the Bernstein basis. Witness calculations certify
existence at those primitives; the sign certificates certify entire domains.
Neither a numerical grid nor a Lean statement is used to certify these
alternative-demand equilibria.

\subsection{Blocking-rule boundary}
At high $F$, a switch from $SU_{12}$ to $SU_{13}$ gives
$W_1^{13,N}=W_1^{12,N}$ and $W_3^{13,N}>W_3^{12,N}$.
It is a weak/Pareto block but not a strict block. The Lean theorem
\path{FV5_indifferenceSeparatesBlockingRules} verifies this logical predicate
distinction, conditional on the payoff comparisons; it does not derive them.
For $m_1=m_2=1$, $m_3=1-\delta$,
\[
W_1^{12,N}-W_1^{13,N}=\delta(A-C)>0.
\]
In the intermediate full-bypass continuation the IS gains are
$m_o(P-C)>0$ for a member and $m_oJ+F>0$ for the outsider, where
$J=K_I+P-K_O-D$. Appendix G enumerates the induced partitions and
the exact asymmetric witness ranges. These statements concern $d=0$;
the positive-$d$ government-ranking result is not an incomplete-adaptation
coalition-stability theorem.
